\documentclass{book}

\title{Science Textbook}
\newcommand{\booksubtitle}{A Latex Template for a Science Textbook}
\newcommand{\booklicense}{Creative Commons Zero 1.0 Universal}

\author{Author Name}
% Author subtitle could be a university or a geographical location, for example
\newcommand{\authorsubtitle}{City, Country}

% Create convenient commands \booktitle and \bookauthor
\makeatletter
\newcommand{\booktitle}{\@title}
\newcommand{\bookauthor}{\@author}
\makeatother

% This utf8 declaration is not needed for versions of latex > 2018 but may
% be helpful for older software. Eventually it may not be worth keeping.
\usepackage[utf8]{inputenc}  
\usepackage{fix-cm}  % this package allows large \fontsize
\usepackage{tikz}    % this is for graphics. e.g. rectangle on title page
\usepackage{amsmath} % Used by equations

% The following dimensions specify 4.75" X 7.5" content on 6 3/8" by 9 1/4"
% paper. The paper width and height can be tweaked as required and the content
% should size to fit within the margins accordingly.
%
% The (inside) bindingoffset should be larger for books with more pages. Some
% standard recommended sizes are .375in minimum up to 1in for 600+ page books.
% Sizes .75in and .875in are also recommended roughly at 150 and 400 pages.
\usepackage[bindingoffset=0.625in,
            left=.5in, right=.5in,
            top=.8125in, bottom=.9375in,
            paperwidth=6.375in, paperheight=9.25in]{geometry}
% Here is an alternative geometry for reading on letter size paper:
% \usepackage[margin=.75in, paperwidth=8.5in, paperheight=11in]{geometry}

\renewcommand{\contentsname}{Table of Contents} % default is {Contents}
\usepackage{makeidx}
\makeindex % Initialize an index so we can add entries with \index

% The next few commands are for creating fake content to fill out the template.
% You should delete this (e.g.  everything up to, but not including,
% \begin{document}) after you insert your own content.
% Example content from Einstein's Meaning of Relativity.
% Public domain book: http://www.gutenberg.org/ebooks/36276
\newcommand{\fakeparagraph}{The theory of relativity is intimately connected
with the theory of space and time. I shall therefore begin with a brief
investigation of the origin of our ideas of space and time, although in
doing so I know that I introduce a controversial subject. The object of all
science, whether natural science or psychology, is to co-ordinate our
experiences and to bring them into a logical system. How are our customary
ideas of space and time related to the character of our experiences?}

\newcommand{\fakecontent}{
\fakeparagraph{}
\fakeparagraph{}

\fakeparagraph{}}

\newcommand{\fakesections}{
\fakecontent{}
\section{First Principles}\fakecontent{}
\subsection{Examples}\fakecontent{}
\subsection{Excessive Elaborations}\fakecontent{}
\subsection{Long Winded Conclusion}\fakecontent{}
\subsection{Exercises}\fakecontent{}
\section{Theory vs Practice}\fakecontent{}
\subsection{Examples}\fakecontent{}
\subsection{Excessive Elaborations}\fakecontent{}
\subsection{Long Winded Conclusion}\fakecontent{}
\subsection{Exercises}\fakecontent{}
\foreach \x in {A,B,C,D,E,F,G,H,I,J,K,L,M,N,O,P,Q,R,S,T,U,V,W,X,Y,Z}
    {\index{\x 1}\index{\x 2}}
}

% Content Starts Here
\begin{document}
\frontmatter

% ---- Half Title Page ----
% current geometry will be restored after title page
\newgeometry{top=1.75in,bottom=.5in}
\begin{titlepage}
\begin{flushleft}

% Title
\textbf{\fontfamily{qcs}\fontsize{48}{54}\selectfont Science\\Textbook\\}

% Draw a line 4pt high
\par\noindent\rule{\textwidth}{4pt}\\

% Subtitle
% Shaded box from left to right. Text node is midway (centered).
\begin{tikzpicture}
\shade[bottom color=lightgray,top color=white]
    (0,0) rectangle (\textwidth, 1.5)
    node[midway] {\textbf{\large \textit{\booksubtitle}}};
\end{tikzpicture}

% Edition Number
\begin{flushright}
\Large First Edition
\end{flushright}

% \vspace{\fill}
\vspace{\fill}

\end{flushleft}
\end{titlepage}
\restoregeometry
% ---- End of Half Title Page ----


% No page numbers on the Frontispiece page
\thispagestyle{empty}


% ---- Title Page ----
% current geometry will be restored after title page
\newgeometry{top=1.75in,bottom=.5in}
\begin{titlepage}
\begin{flushleft}

% Title
\textbf{\fontfamily{qcs}\fontsize{48}{54}\selectfont Science\\Textbook\\}

% Draw a line 4pt high
\par\noindent\rule{\textwidth}{4pt}\\

% Shaded box from left to right with Subtitle
% The text node is midway (centered).
\begin{tikzpicture}
\shade[bottom color=lightgray,top color=white]
    (0,0) rectangle (\textwidth, 1.5)
    node[midway] {\textbf{\large \textit{\booksubtitle}}};
\end{tikzpicture}

% Edition Number
\begin{flushright}
\Large First Edition
\end{flushright}

\vspace{\fill}

% Author and Location
\textbf{\large \bookauthor}\\[3.5pt]
\textbf{\large \textit{\authorsubtitle}}

\vspace{\fill}

% Self Publishing Logo. Free to use: CC0 license.
% The source file is book.svg. If you change the svg, you must then convert
% it to pdf. There are many online and offline tools available to do that.
\begin{center}
\includegraphics{booksvg.pdf}\\[4pt]
\fontfamily{lmtt}\small{Self Publishers Worldwide\\
Seattle San Francisco New York\\
London Paris Rome Beijing Barcelona}
\end{center}

\end{flushleft}
\end{titlepage}
\restoregeometry
% ---- End of Title Page ----

% Do not show page numbers on colophon page
\thispagestyle{empty}

\begin{flushleft}
\vspace*{\fill}
This book was typeset using \LaTeX{} software.\\
\vspace{\fill}
Copyright \textcopyright{} \the\year{}  \bookauthor\\
License: \booklicense
\end{flushleft}

% A title page resets the page # to 1, but the second title page
% was actually page 3. So add two to page counter.
\addtocounter{page}{2}

% The asterisk excludes chapter from the table of contents.
\chapter*{Preface}
\fakeparagraph{}
\[
\left.
\begin{aligned}
  &\sqrt{{dX_{1}}^{2} + {dX_{2}}^{2} + {dX_{3}}^{2}} \\
  &\qquad
  \begin{aligned}
  &= \left(1 + \frac{\kappa}{8\pi} \int \frac{\sigma\, dV_{0}}{r}\right)
     \sqrt{{dx_{1}}^{2} + {dx_{2}}^{2} + {dx_{3}}^{2}}, \\
dT &= \left(1 - \frac{\kappa}{8\pi} \int \frac{\sigma\, dV_{0}}{r}\right) dl.
\end{aligned}
\end{aligned}
\right\}
\]
\fakecontent{}
\fakecontent{}

% Three-level Table of Contents
\setcounter{tocdepth}{3}
\tableofcontents

\mainmatter

\chapter{Introduction}\fakesections{}
\chapter{Theory of Numbers}\fakesections{}
\chapter{Irrational and Transcendent Numbers}\fakesections{}
\chapter{Complex Numbers}\fakesections{}
\chapter{Quaternions and Ausdehnungslehre}\fakesections{}
\chapter{Theory of Equations}\fakesections{}
\chapter{Substitutions and Groups}\fakesections{}
\chapter{Determinants}\fakesections{}
\chapter{Quantics}\fakesections{}
\chapter{Calculus}\fakesections{}
\chapter{Differential Equations}\fakesections{}
\chapter{Infinite Series}\fakesections{}
\chapter{Theory of Functions}\fakesections{}
\chapter{Probabilities and Least Squares}\fakesections{}
\chapter{Analytic Geometry}\fakesections{}
\chapter{Modern Geometry}\fakesections{}
\chapter{Trigonometry and Elementary Geometry}\fakesections{}
\chapter{Non-Euclidean Geometry}\fakesections{}
\chapter{Bibliography}\fakesections{}
\chapter{General Tendencies}\fakesections{}

\backmatter
\addcontentsline{toc}{chapter}{Index}
\printindex
\end{document}